% =====================================================================
%  Synopsis Report  —  "KioMeet: A Secure Native Interview & Conferencing Platform"
%  Formatting per guidelines:
%    - Printed one side only
%    - 1.5 line spacing
%    - Margins: left 2.5cm, top 2.5cm, right 1.25cm, bottom 1.25cm
% =====================================================================
% NOTE: compile with XeLaTeX (uses the real macOS system fonts):
%   xelatex synopsis.tex
\documentclass[12pt,oneside,a4paper]{article}

% ---- Page geometry (margins per guideline) --------------------------
\usepackage[a4paper,left=2.5cm,top=2.5cm,right=1.25cm,bottom=1.25cm]{geometry}

% ---- Core packages --------------------------------------------------
\usepackage{setspace}        % line spacing
\usepackage{graphicx}        % figures / logo
\usepackage{array}           % table column tuning
\usepackage{xcolor}
\usepackage[hidelinks]{hyperref}

% Name the auto-generated bibliography heading to match the other sections.
\renewcommand{\refname}{REFERENCES}

% ---- Fonts (match the reference report) -----------------------------
% The sample report's body is an Arial/Helvetica-style sans-serif, with the
% main title and the degree block in a Times-style serif. Using fontspec we
% load the real system fonts: Arial for the body, Times New Roman for the
% serif title-page lines (via \titleserif).
\usepackage{fontspec}
\setmainfont{Arial}
\newfontfamily\titleserif{Times New Roman}

% ---- 1.5 line spacing ----------------------------------------------
\onehalfspacing

% ---- Paragraph spacing (parskip-style) -----------------------------
\setlength{\parindent}{0pt}
\setlength{\parskip}{6pt plus 2pt minus 1pt}

% Note: \section* in the article class already renders as a bold,
% large, upright heading, which matches the sample report.

\begin{document}

% =====================================================================
%  TITLE PAGE
% =====================================================================
\begin{titlepage}
\begin{singlespace}
\centering
\vspace*{0.3cm}
{\titleserif\Huge\bfseries Synopsis Report\par}
\vspace{1.0cm}

{\large\itshape ``KioMeet --- A Secure Native Interview \&\\[2pt]
Video Conferencing Platform''\par}
\vspace{0.4cm}

{\large\itshape by\par}
\vspace{0.25cm}
{\large\itshape Aman Verma\par}
{\large\itshape 23bcs012\par}
\vspace{0.1cm}
{\large\itshape Ankit Kashyap\par}
{\large\itshape 23bcs017\par}
\vspace{0.1cm}
{\large\itshape Amrit Kashyap\par}
{\large\itshape 23bcs014\par}

\vspace{0.8cm}
{\large\itshape\underline{Under the guidance of Dr. Manoj Gupta}\par}

\vspace{0.8cm}
{\titleserif\large Submitted in partial fulfilment of the requirements for the\\
award of the degree\par}

\vspace{0.7cm}
{\titleserif\large BACHELOR OF TECHNOLOGY\par}
{\titleserif\large in\par}
{\titleserif\large COMPUTER SCIENCE \& ENGINEERING\par}

\vspace{0.7cm}
\IfFileExists{logo.png}{\includegraphics[width=3.0cm]{logo.png}\par}{}

\vfill
{\large SHRI MATA VAISHNO DEVI UNIVERSITY, KATRA\par}
{\large (School of Computer Science \& Engineering)\par}
{\large JAMMU \& KASHMIR -- 182 320\par}
{\large Session 2026--27 (Odd)\par}
\vspace{0.4cm}
\end{singlespace}
\end{titlepage}

% =====================================================================
%  INTRODUCTION
% =====================================================================
\section*{INTRODUCTION}
\addcontentsline{toc}{section}{Introduction}

Remote hiring and distributed teamwork have made video conferencing a daily
necessity. Yet most general-purpose meeting tools are browser- or
Electron-based, which limits access to native operating-system capabilities,
adds overhead, and --- critically for technical hiring --- offers no reliable
way to stop private interviewer material from leaking into the shared video
stream.

This project proposes \textbf{KioMeet}, a native cross-platform video
conferencing application for \textbf{Windows and macOS}, built specifically for
teams and technical interviews. KioMeet runs as a native desktop application on
each platform (Objective-C and Swift on macOS; C++/.NET on Windows) with no
browser or Electron layer, pairing a GPU-accelerated media pipeline with a
dedicated Node.js backend for room management, authentication, scheduling,
recording, and team collaboration.

KioMeet's distinguishing feature is a dedicated \textbf{secure interview mode}.
In this mode the interviewer's own screen is rendered inside a hardened window
whose contents are never captured into the outgoing stream; only the content
of designated tool surfaces (code editor, whiteboard) is shared with the
candidate. This guarantees that private questions, notes, browser tabs, or
other application windows never bleed into the call, preserving interview
integrity while keeping the experience fast and natural.

\subsection*{System Overview}
KioMeet is organised in three cooperating layers. The \emph{native desktop
client} (Windows and macOS) owns media capture, GPU rendering (Metal on macOS,
Direct3D on Windows), the secure window, and the user interface. The
\emph{real-time media layer} uses a LiveKit-compatible WebRTC
server~\cite{webrtc,livekit} to publish and subscribe audio, video, and
screen-share tracks. The
\emph{backend token server} (Node.js + Express) issues short-lived access
tokens, manages rooms, enforces user tiers, schedules meetings, syncs
calendars, and tracks recordings, backed by PostgreSQL for persistent data and
Redis for ephemeral room state.

% =====================================================================
%  LITERATURE SURVEY
% =====================================================================
\section*{LITERATURE SURVEY}
\addcontentsline{toc}{section}{Literature Survey}

\subsection*{1. WebRTC and the LiveKit SFU model}
WebRTC is the open standard for browser- and application-based real-time
communication~\cite{webrtc,rfc8825}. At scale, mesh topologies become
impractical, so modern platforms adopt a Selective Forwarding Unit (SFU) such
as LiveKit~\cite{livekit}, which routes media without transcoding and supports
adaptive simulcast. This body of work
shows that an SFU is the appropriate architecture for rooms of up to tens of
participants while keeping client CPU and bandwidth bounded.

\subsection*{2. Native screen-capture APIs (Windows \& macOS)}
Both desktop platforms expose modern, permission-gated screen-capture APIs:
Apple's ScreenCaptureKit on macOS~\cite{screencapturekit} and the Windows
Graphics Capture API on Windows~\cite{wgc}, each supporting display and
individual-window capture with live thumbnails and system-audio capture. Compared with legacy capture paths they
offer lower latency, finer-grained window selection, and explicit user consent.
This establishes that native frameworks on each OS can deliver screen sharing
that is both higher quality and more privacy-respecting than cross-platform
wrappers.

\subsection*{3. Secure window composition and content isolation}
Research and platform documentation on window-server
composition~\cite{screencapturekit,wgc} show that an application can place a
window above the normal layer order and selectively exclude surfaces from
capture. Applied to interviewing, this means the
interviewer's private material can be isolated from the shared stream at the
compositor level, so only explicitly designated tool surfaces are ever
transmitted --- a stronger guarantee than hiding windows in software alone.

\subsection*{4. Token-based authentication and tiered access}
Established security practice favours a short-lived JWT access
token~\cite{jwt} paired with a rotating, hashed refresh token, with the user's
service tier re-read from the
database on every refresh rather than trusted from the client. This literature
underpins KioMeet's auth model and its fail-closed, event-driven billing and tier
enforcement.

\subsection*{SUMMARY OF LITERATURE SURVEY}
\begin{enumerate}
  \item A native client built on an SFU (LiveKit) and each platform's modern
        frameworks (ScreenCaptureKit/Metal on macOS, Windows Graphics
        Capture/Direct3D on Windows) delivers lower-latency, higher-quality
        conferencing than browser or Electron alternatives.
  \item Content isolation for interviews is best enforced at the window-server
        / capture-path level, not by hiding windows in application code.
  \item A short-lived access token plus rotating refresh token, with tier always
        re-read server-side, is the accepted pattern for secure, tier-gated
        access.
\end{enumerate}

% =====================================================================
%  PROBLEM STATEMENT
% =====================================================================
\section*{PROBLEM STATEMENT}
\addcontentsline{toc}{section}{Problem Statement}

To design and build a native cross-platform (Windows and macOS) video
conferencing platform that supports both standard team calls and technical
interviews, in which the interviewer's
private material is provably isolated from the shared video stream, while
delivering reliable multi-participant audio/video, screen sharing, recording,
scheduling, and moderation through a secure, tier-gated backend --- without the
overhead and capability limits of browser- or Electron-based solutions.

% =====================================================================
%  OBJECTIVES
% =====================================================================
\section*{OBJECTIVES}
\addcontentsline{toc}{section}{Objectives}

The specific objectives of this project are as follows:
\begin{enumerate}
  \item To build native desktop clients for \textbf{Windows and macOS}
        (Objective-C/Swift on macOS; C++/.NET on Windows) that capture camera,
        microphone, and screen content and render remote participants on the
        GPU (Metal on macOS, Direct3D on Windows), scaling the layout adaptively
        up to 50 participants.
  \item To implement a \emph{secure interview mode} in which the interviewer's
        screen is composed into a hardened window and only designated tool
        surfaces (code editor, whiteboard) are streamed to the candidate,
        preventing any content leakage.
  \item To integrate a LiveKit-compatible WebRTC server for publishing and
        subscribing audio, video, and screen-share tracks, with decoupled
        audio/video paths so a microphone toggle never interrupts the camera.
  \item To develop a Node.js backend that issues short-lived JWT access tokens
        with rotating refresh tokens, manages rooms (Redis, 24-hour TTL), and
        stores persistent data in PostgreSQL.
  \item To provide meeting features --- public and direct chat, raise-hand and
        speaker queue, live polls, a Q\&A board, role-based moderation, a
        lobby/waiting room, and local, cloud, and multi-track recording.
  \item To support scheduling with email invitations (\texttt{.ics}) and
        Apple, Google, and Outlook calendar sync, and tier-gated access (free,
        pro, hiring) that is enforced server-side and fail-closed.
\end{enumerate}

% =====================================================================
%  METHODOLOGY
% =====================================================================
\section*{METHODOLOGY}
\addcontentsline{toc}{section}{Methodology}

\subsection*{1. Requirement Analysis}
\begin{itemize}
  \item Identify needs for team calls and technical interviews (isolation,
        recording, moderation, scheduling).
  \item Study existing platforms (Zoom, Google Meet, native WebRTC apps) and
        the LiveKit SFU model.
\end{itemize}

\subsection*{2. System Design}
\begin{itemize}
  \item \textbf{Native clients (Windows \& macOS)} --- media capture, GPU
        rendering (Metal/Direct3D), secure window, and UI panels.
  \item \textbf{Media layer} --- LiveKit publisher/subscriber bridging local
        and remote tracks.
  \item \textbf{Backend} --- Express token server with PostgreSQL (persistent)
        and Redis (ephemeral room state).
\end{itemize}

\subsection*{3. Core Client Development (Windows \& macOS)}
\begin{itemize}
  \item Camera and microphone capture with two-phase toggles that eliminate
        frozen/black frames on mute.
  \item Screen and window sharing via each platform's native capture API
        (ScreenCaptureKit on macOS, Windows Graphics Capture on Windows) with a
        live-thumbnail picker and optional system-audio capture.
  \item Secure interview mode: interviewer content composed into a hardened
        window; only designated tool surfaces are published to the candidate.
  \item Adaptive grid, stage + filmstrip, and spotlight layouts; active-speaker
        highlighting and auto-spotlight.
\end{itemize}

\subsection*{4. Backend \& Real-Time Services}
\begin{itemize}
  \item JWT access tokens (short TTL) with rotating, hashed refresh tokens and
        family-based theft detection.
  \item Room creation and 6-character join codes with 24-hour expiry in Redis.
  \item Scheduling with \texttt{.ics} email invites and Google/Outlook OAuth
        calendar sync; team management with role-based permissions.
  \item Local, cloud (S3 egress), and multi-track recording with a composed
        renderer and disk-space safeguards.
\end{itemize}

\subsection*{5. Testing, Deployment \& Verification}
\begin{itemize}
  \item Unit and integration tests on each platform (XCTest on macOS and the
        equivalent on Windows) covering tokens, media bridges, room lifecycle,
        isolation invariants, and recording.
  \item Health-checked backend deployment (PostgreSQL + Redis) and end-to-end
        verification between live Windows and macOS clients.
\end{itemize}

% =====================================================================
%  EXPECTED RESULTS
% =====================================================================
\section*{EXPECTED RESULTS}
\addcontentsline{toc}{section}{Expected Results}

\begin{enumerate}
  \item \textbf{Reliable native conferencing} --- smooth multi-participant
        audio/video with GPU rendering and adaptive layouts, outperforming
        browser/Electron alternatives in latency and resource use.
  \item \textbf{Guaranteed interview isolation} --- in secure interview mode,
        only designated tool surfaces reach the candidate; the interviewer's
        other windows and material never appear in the stream.
  \item \textbf{Full collaboration suite} --- chat, polls, Q\&A, raise-hand
        queue, moderation, lobby, scheduling with calendar sync, teams, and
        local/cloud/multi-track recording working end-to-end.
  \item \textbf{Secure, tier-gated backend} --- short-lived tokens, rotating
        refresh tokens, and fail-closed tier enforcement protecting all
        features and data.
\end{enumerate}

% =====================================================================
%  TIME SCHEDULE
% =====================================================================
\section*{TIME SCHEDULE}
\addcontentsline{toc}{section}{Time Schedule}

\renewcommand{\arraystretch}{1.5}
\noindent
\begin{tabular}{|p{5.6cm}|p{2.6cm}|p{4.6cm}|}
\hline
\textbf{Phase} & \textbf{Duration} & \textbf{Timeline} \\ \hline
Phase 1: Literature Review \& Requirements & 2 weeks & 1 Sept -- 14 Sept 2026 \\ \hline
Phase 2: System Design & 2 weeks & 15 Sept -- 28 Sept 2026 \\ \hline
Phase 3: Client \& Backend Implementation & 4 weeks & 29 Sept -- 26 Oct 2026 \\ \hline
Phase 4: Integration (media, secure mode) & 2 weeks & 27 Oct -- 9 Nov 2026 \\ \hline
Phase 5: Testing \& Debugging & 2 weeks & 10 Nov -- 23 Nov 2026 \\ \hline
Phase 6: Documentation \& Finalization & 1 week & 24 Nov -- 30 Nov 2026 \\ \hline
\textbf{Total Project Duration} & \textbf{13 weeks} & \textbf{1 Sept -- 30 Nov 2026} \\ \hline
\end{tabular}

% =====================================================================
%  REFERENCES
% =====================================================================
\addcontentsline{toc}{section}{References}
\begin{thebibliography}{9}
  \bibitem{webrtc} WebRTC --- Official Project Documentation. \url{https://webrtc.org}
  \bibitem{livekit} LiveKit --- Open-Source WebRTC SFU Documentation. \url{https://docs.livekit.io}
  \bibitem{screencapturekit} Apple Developer Documentation --- ScreenCaptureKit \& Metal (macOS). \url{https://developer.apple.com/documentation/screencapturekit}
  \bibitem{wgc} Microsoft Documentation --- Windows Graphics Capture \& Direct3D (Windows). \url{https://learn.microsoft.com/windows/uwp/audio-video-camera/screen-capture}
  \bibitem{rfc8825} IETF RFC 8825 --- Overview: Real-Time Protocols for Browser-Based Applications (WebRTC).
  \bibitem{jwt} IETF RFC 7519 --- JSON Web Token (JWT).
\end{thebibliography}

% =====================================================================
%  SIGNATURES
% =====================================================================
\vspace{1.4cm}
\noindent
\begin{minipage}[t]{0.5\textwidth}
\textbf{Student Sign}
\begin{enumerate}
  \item Aman Verma
  \item Ankit Kashyap
  \item Amrit Kashyap
\end{enumerate}
\end{minipage}%
\begin{minipage}[t]{0.5\textwidth}
\raggedleft
\textbf{Guide Sign}
\end{minipage}

\end{document}
